\documentclass[11pt,letterpaper]{article}

\usepackage[utf8]{inputenc}
\usepackage[T1]{fontenc}
\usepackage[spanish,es-noquoting]{babel}
\usepackage[margin=2.5cm]{geometry}
\usepackage{amsmath}
\usepackage{listings}
\usepackage{xcolor}
\usepackage{enumitem}
\usepackage[hidelinks]{hyperref}

\definecolor{grisfondo}{gray}{0.96}

\lstset{
  basicstyle=\ttfamily\small,
  backgroundcolor=\color{grisfondo},
  frame=single,
  framesep=5pt,
  breaklines=true,
  columns=fullflexible,
  showstringspaces=false
}

\title{\textbf{Derivadas simbólicas en Prolog}}
\author{Práctica de Lógica Computacional}
\date{\today}

\begin{document}

\maketitle

\begin{center}
\fbox{\parbox{0.85\textwidth}{\centering
Código fuente del proyecto:

\url{https://github.com/hmndz3/LogicaPY3}
}}
\end{center}

\vspace{1em}

% =====================================================================
\section{¿Qué pedía el enunciado?}

Nos pidieron escribir un programa en Prolog que calcule derivadas. No la
derivada evaluada en un número, sino la derivada \emph{escrita como fórmula}:
uno le da una expresión matemática y el programa devuelve la expresión de su
derivada.

Las funciones que tenía que soportar son:

\begin{itemize}[noitemsep]
  \item Suma, resta, multiplicación y división.
  \item Funciones polinomiales, es decir potencias como $x^5$.
  \item Seno, coseno y tangente.
  \item Arcotangente.
  \item Exponencial y logaritmo natural.
\end{itemize}

Además tenía que aplicar la \textbf{regla de la cadena}: cuando una función
está metida dentro de otra, como $\mathrm{sen}(e^{3x})$, hay que derivar la de
afuera y multiplicar por la derivada de la de adentro.

Hay una condición importante: \textbf{el resultado no se debe simplificar}.
Tiene que quedar indicado tal cual sale de aplicar las reglas. Por eso en los
ejemplos del enunciado aparecen cosas como \texttt{1+0} o \texttt{2*x\^{}1} en
lugar de \texttt{1} o \texttt{2*x}. El programa no hace álgebra, solo aplica
reglas de derivación.

% =====================================================================
\section{¿Qué hicimos?}

Escribimos un solo archivo, \texttt{derivadas.pl}, con un predicado llamado
\texttt{derivada} que recibe tres cosas:

\begin{lstlisting}
derivada(Expresion, Variable, Resultado)
\end{lstlisting}

Uno le pasa la expresión y la variable respecto a la que quiere derivar, y
Prolog le devuelve el resultado.

\subsection{La idea de fondo}

Prolog no funciona como los lenguajes a los que estamos acostumbrados. Uno no
le escribe pasos, le escribe \textbf{reglas}. En este caso le escribimos, una
por una, las mismas reglas de derivación que vimos en Cálculo:

\begin{center}
\renewcommand{\arraystretch}{1.8}
\begin{tabular}{ll}
\textbf{Regla de Cálculo} & \textbf{Cómo quedó en el programa} \\
$(u+v)' = u' + v'$ & \texttt{derivada(U+V, X, DU+DV)} \\
$(u \cdot v)' = u'v + uv'$ & \texttt{derivada(U*V, X, DU*V+U*DV)} \\
$\left(\dfrac{u}{v}\right)' = \dfrac{u'v - uv'}{v^2}$ & \texttt{derivada(U/V, X, (DU*V-U*DV)/V\^{}2)} \\
$(\mathrm{sen}\,u)' = \cos(u) \cdot u'$ & \texttt{derivada(sen(U), X, cos(U)*DU)} \\
$(\ln u)' = \dfrac{u'}{u}$ & \texttt{derivada(ln(U), X, DU/U)} \\
\end{tabular}
\end{center}

Lo que hace Prolog es ir buscando cuál de esas reglas encaja con la expresión
que le dimos. Si le damos una suma, usa la regla de la suma. Si le damos un
logaritmo, usa la del logaritmo. Y como cada regla vuelve a llamar a
\texttt{derivada} sobre las partes de adentro, la regla de la cadena sale sola:
no hubo que programarla aparte.

\subsection{Un ejemplo paso a paso}

Si le pedimos la derivada de $\mathrm{sen}(e^{3x})$, el programa hace esto:

\begin{enumerate}[noitemsep]
  \item Ve un seno, así que aplica $\cos(\text{adentro}) \times (\text{derivada de adentro})$.
  \item Para derivar lo de adentro, $e^{3x}$, ve una exponencial y aplica
        $e^{3x} \times (\text{derivada de } 3x)$.
  \item Para derivar $3x$ ve un número multiplicando, así que lo deja quieto y
        deriva la $x$, que da $1$.
\end{enumerate}

Juntando todo:

\begin{lstlisting}
?- derivada(sen(exp(3*x)), x, Y).
Y = cos(exp(3*x))*(exp(3*x)*(3*1)).
\end{lstlisting}

Que es exactamente lo que pedía el enunciado.

\subsection{Organización del código}

El archivo está dividido en bloques, uno por cada regla, separados con un
comentario que dice qué hace cada uno:

\begin{lstlisting}
% ------------------------------------------------ division
% d(u/v) = (du*v - u*dv) / v^2
derivada(U/V, X, (DU*V-U*DV)/V^2) :- !, derivada(U,X,DU), derivada(V,X,DV).
\end{lstlisting}

Al final del archivo hay un bloque de pruebas. Ahí están guardados los cinco
ejemplos del enunciado junto con el resultado que deberían dar, más cinco
ejemplos nuestros para cubrir coseno, tangente, resta, producto y exponencial.

% =====================================================================
\section{Resultados}

Corrimos las pruebas y las diez salen correctas. Los cinco ejemplos del
enunciado dan el mismo resultado, carácter por carácter:

\begin{lstlisting}
OK    derivada(ln(x/(x+1)),x) = (1*(x+1)-x*(1+0))/(x+1)^2/(x/(x+1))
OK    derivada(sen(exp(3*x)),x) = cos(exp(3*x))*(exp(3*x)*(3*1))
OK    derivada(ln(sen((x+1)/x)),x) = cos((x+1)/x)*(((1+0)*x-(x+1)*1)/x^2)
                                     /sen((x+1)/x)
OK    derivada((x^5+x^3)/(x^2+1),x) = ((5*x^4+3*x^2)*(x^2+1)
                                       -(x^5+x^3)*(2*x^1+0))/(x^2+1)^2
OK    derivada(arctan((x^2+1)^10),x) = 10*(x^2+1)^9*(2*x^1+0)
                                       /(1+((x^2+1)^10)^2)
OK    derivada(cos(x^2),x) = -sen(x^2)*(2*x^1)
OK    derivada(tan(x),x) = (1+tan(x)^2)*1
OK    derivada(x^3-x,x) = 3*x^2-1
OK    derivada(x*sen(x),x) = 1*sen(x)+x*(cos(x)*1)
OK    derivada(exp(x)/x,x) = (exp(x)*1*x-exp(x)*1)/x^2

10 de 10 pruebas correctas
\end{lstlisting}

% =====================================================================
\section{Cómo correrlo}

Se necesita SWI-Prolog instalado. Desde la terminal, parado en la carpeta del
proyecto:

\begin{lstlisting}
swipl derivadas.pl
\end{lstlisting}

Eso abre la consola de Prolog. Ahí se pueden correr las pruebas:

\begin{lstlisting}
?- pruebas.
\end{lstlisting}

O pedir una derivada cualquiera:

\begin{lstlisting}
?- derivada(ln(x/(x+1)), x, Y).
\end{lstlisting}

Para salir se escribe \texttt{halt.}

% =====================================================================
\section{Referencias}

\begin{itemize}[noitemsep]
  \item \url{http://www.nivaria.net/blog/index.php/prolog/383-derivada-de-una-funcion-entera}
  \item \url{https://gist.github.com/tokkenno/8038482}
  \item \url{http://turing.iimas.unam.mx/~pedroG/al-2008-II/Lab_05.pdf}
  \item Manual de SWI-Prolog: \url{https://www.swi-prolog.org/pldoc/doc_for?object=manual}
\end{itemize}

\end{document}
